% =============================================================================
% ForensiQ — Mapeamento de Requisitos ISO/IEC 27037 (v1.2)
% UC 21184 · Projeto de Engenharia Informática · UAb
% Versão 1.2 (3 maio 2026): Refeita a partir do código real para corrigir
% desvios da v1.1 (12 abr 2026), que referia estados em inglês inexistentes
% (RECEIVED, IN_EXAMINATION, ...) e tipos antigos (USB_DRIVE, HARD_DRIVE,
% SD_CARD) substituídos pela taxonomia digital-first do ADR-0010.
% =============================================================================
\documentclass[a4paper,11pt]{article}

\usepackage[utf8]{inputenc}
\usepackage[T1]{fontenc}
\usepackage[top=2.5cm, bottom=2.5cm, left=3cm, right=2.5cm]{geometry}
\usepackage{booktabs}
\usepackage{longtable}
\usepackage{array}
\usepackage{parskip}
\usepackage{hyperref}

\hypersetup{
    colorlinks=true,
    linkcolor=black,
    urlcolor=blue,
    pdfauthor={João Rodrigues},
    pdftitle={ForensiQ — Mapeamento ISO/IEC 27037 v1.2}
}

\newcolumntype{L}[1]{>{\raggedright\arraybackslash}p{#1}}

% =============================================================================
\begin{document}
% =============================================================================

\begin{center}
    {\LARGE\bfseries Mapeamento de Requisitos ISO/IEC 27037}\\[0.4em]
    {\large ForensiQ --- Plataforma Modular de Gestão de Prova Digital\\
    para \emph{First Responders}}\\[1.2em]
    \begin{tabular}{ll}
        \textbf{Versão:} & 1.2 · 3 maio 2026 \\
        \textbf{Norma:}  & ISO/IEC 27037:2012 \\
        \textbf{Estado:} & Aceite, alinhada ao código de produção \\
    \end{tabular}
\end{center}

\vspace{0.5em}\hrule\vspace{1.5em}

\section{Introdução}

Este documento apresenta a matriz de rastreabilidade entre os requisitos
da norma ISO/IEC~27037:2012 (\emph{Guidelines for identification,
collection, acquisition and preservation of digital evidence}) e as
funcionalidades implementadas na plataforma ForensiQ. A versão~1.2
substitui a~1.1 (12~abr~2026), que descrevia estados em inglês e tipos
de evidência inexistentes no código actual; a presente versão é
refeita a partir do código de produção e ancora cada controlo num
caminho \texttt{ficheiro:linha} verificável.

\section{Matriz de Rastreabilidade}

\subsection*{§5.1 --- Princípios gerais}

\textbf{Implementação.} Controlo de acesso baseado em papéis com dois
perfis operacionais: \texttt{AGENT} (recolha de prova) e \texttt{EXPERT}
(perícia), distinguidos pelo campo \texttt{User.profile}
(\texttt{src/backend/core/models.py:123--129}). O perfil
administrativo é representado pela flag \texttt{is\_staff} do Django.
\emph{Tokens} JWT em \emph{cookies} HttpOnly~+ Secure~+ SameSite=Strict
(\texttt{src/backend/core/auth.py:33--58}), com rotação de \emph{refresh}
e \emph{blacklist}
(\texttt{src/backend/forensiq\_project/settings.py:174--184}). Auditoria
transversal em \texttt{AuditLog} com \texttt{correlation\_id} por pedido
(\texttt{models.py:1199--1206}; \texttt{audit.py:83--134}).

\subsection*{§5.2 --- Identificação de prova digital}

\textbf{Implementação.} Modelo \texttt{Evidence}
(\texttt{models.py:306--724}) com campos obrigatórios:
\begin{itemize}
    \item \texttt{type} (\texttt{EvidenceType} \emph{enum},
          \texttt{models.py:332--353}) com 18 tipos digital-first
          (ADR-0010): 14 raízes (\texttt{MOBILE\_DEVICE},
          \texttt{COMPUTER}, \texttt{STORAGE\_MEDIA},
          \texttt{GAMING\_CONSOLE}, \texttt{GPS\_TRACKER},
          \texttt{SMART\_TAG}, \texttt{CCTV\_DEVICE},
          \texttt{VEHICLE}, \texttt{DRONE}, \texttt{IOT\_DEVICE},
          \texttt{NETWORK\_DEVICE}, \texttt{DIGITAL\_FILE},
          \texttt{RFID\_NFC\_CARD}, \texttt{OTHER\_DIGITAL}) e 4
          sub-componentes típicos (\texttt{SIM\_CARD},
          \texttt{MEMORY\_CARD}, \texttt{INTERNAL\_DRIVE},
          \texttt{VEHICLE\_COMPONENT});
    \item \texttt{description}, \texttt{serial\_number},
          \texttt{gps\_lat}/\texttt{gps\_lon}, \texttt{timestamp\_seizure}
          --- todos validados em \texttt{Evidence.clean}
          (\texttt{models.py:630--723});
    \item identificador legível \texttt{code} no formato
          \texttt{ITM-YYYY-NNNNN}, gerado dentro de transacção atómica com
          retentativa (\texttt{\_next\_yearly\_code},
          \texttt{models.py:35--85}).
\end{itemize}

\subsection*{§5.3 --- Recolha e aquisição de prova}

\textbf{Implementação.} Registo fotográfico obrigatório em
\texttt{Evidence.photo} (\texttt{ImageField} validado por
\texttt{validate\_image\_max\_size}, máximo 25~MB).
Geo-referenciação capturada via \texttt{Geolocation API} no
\emph{wizard} \emph{frontend} (\texttt{evidences\_new.html:40--289}, 8
passos) e gravada como \texttt{gps\_lat}/\texttt{gps\_lon}.
\emph{Timestamp} sempre do servidor (\texttt{timezone.now()}, NTP-synced),
nunca do cliente (\texttt{ChainOfCustody.save}, \texttt{models.py:1081}).
Exportação PDF da ocorrência e do item em
\texttt{src/backend/core/pdf\_export.py} (gerador ReportLab) com
declaração ISO/IEC~27037 incorporada.

\subsection*{§5.4 --- Preservação de integridade}

\textbf{Implementação.} \emph{Hash} SHA-256 calculado automaticamente:
\begin{itemize}
    \item Em \texttt{Evidence.save}, a função
          \texttt{compute\_integrity\_hash}
          (\texttt{models.py:554--584}) cobre os metadados e o
          \emph{hash} dos \emph{bytes} da fotografia
          (\texttt{\_compute\_photo\_hash}, \texttt{models.py:534--552});
    \item Em \texttt{ChainOfCustody.save}, a função
          \texttt{compute\_record\_hash}
          (\texttt{models.py:989--1031}) é pura, recebe o
          \emph{previous\_hash} como parâmetro e produz um \emph{hash}
          encadeado, reproduzível por qualquer perito independente.
\end{itemize}

A imutabilidade está implementada em \emph{três camadas} de defesa em
profundidade:
\begin{enumerate}
    \item ORM Django: \texttt{Evidence.save}
          (\texttt{models.py:586--617}) e
          \texttt{ChainOfCustody.save}
          (\texttt{models.py:1044--1107}) rejeitam alterações com
          \texttt{ValidationError}; os respectivos \texttt{delete}
          também (\texttt{models.py:619--624, 1109--1114}).
    \item Django REST Framework: \texttt{EvidenceViewSet}
          (\texttt{views.py:384}), \texttt{DigitalDeviceViewSet}
          (\texttt{views.py:566}) e \texttt{ChainOfCustodyViewSet}
          (\texttt{views.py:602}) declaram
          \texttt{http\_method\_names = ['get', 'post', 'head', 'options']}.
    \item PostgreSQL: \emph{triggers} \texttt{BEFORE UPDATE} e
          \texttt{BEFORE DELETE} sobre \texttt{core\_evidence},
          \texttt{core\_chainofcustody} e \texttt{core\_digitaldevice}
          (\texttt{src/backend/core/migrations/0002\_add\_immutability\_triggers.py:36--115}).
\end{enumerate}

\subsection*{§5.5 --- Cadeia de custódia}

\textbf{Implementação.} Modelo \texttt{ChainOfCustody}
(\texttt{models.py:858--1114}) com registo \emph{append-only}.
Máquina de estados linear de \textbf{7 estados} em português:
\texttt{APREENDIDA} $\rightarrow$ \texttt{EM\_TRANSPORTE} $\rightarrow$
\texttt{RECEBIDA\_LABORATORIO} $\rightarrow$ \texttt{EM\_PERICIA}
$\rightarrow$ \texttt{CONCLUIDA} $\rightarrow$ \texttt{DEVOLVIDA} ou
\texttt{DESTRUIDA}, definida em \texttt{VALID\_TRANSITIONS}
(\texttt{models.py:883--892}). A ordem canónica de cada cadeia é dada
pelo campo \texttt{sequence} e pelo \texttt{UniqueConstraint(evidence,
sequence)} (\texttt{models.py:957--962}). A escrita serializa-se com
\texttt{transaction.atomic} + \texttt{select\_for\_update}
(\texttt{models.py:1064--1097}). O encadeamento de \emph{hashes}
permite, por recálculo independente, detectar qualquer adulteração
posterior.

\subsection*{§6.1 --- Competência do pessoal}

\textbf{Implementação.} Sistema de papéis diferenciados aplicado por
\emph{permission classes} dedicadas em
\texttt{src/backend/core/permissions.py}: \texttt{IsAgent},
\texttt{IsExpert}, \texttt{IsAgentOrExpert} e
\texttt{IsOwnerOrReadOnly}. \texttt{AGENT} cria ocorrências e
evidências e só vê os recursos próprios; \texttt{EXPERT} consulta e
participa nas transições de custódia; \emph{staff} acede ao
\emph{Django Admin} (com prefixo configurável via
\texttt{ADMIN\_URL\_PREFIX}). A filtragem por
\texttt{request.user} aplica-se ao nível do \texttt{get\_queryset()}
de cada \emph{viewset}.

\subsection*{§6.2 --- Documentação do processo}

\textbf{Implementação.} Geração automática de PDF forense em
\texttt{pdf\_export.py}: \texttt{generate\_evidence\_pdf} e
\texttt{generate\_occurrence\_pdf}, com identificação da ocorrência
(NUIPC, GPS, agente), inventário da prova (com indentação para
sub-componentes), cadeia de custódia em tabela cronológica e
declaração de integridade no rodapé com referência explícita à
ISO/IEC~27037:2012. Exportação CSV em \emph{streaming} disponível em
\texttt{/api/occurrences/csv/}, \texttt{/api/evidences/csv/} e
\texttt{/api/custody/csv/} (\texttt{views.py}).

\subsection*{§7.1 --- Dispositivos digitais}

\textbf{Implementação.} Modelo \texttt{DigitalDevice}
(\texttt{models.py:730--851}) regista, em ficha técnica complementar a
\texttt{Evidence}: tipo, marca, modelo, nome comercial,
\texttt{condition}, IMEI (validado por algoritmo de Luhn,
\texttt{validators.py:47--66}), número de série e notas. O
modelo coexiste com a hierarquia
\texttt{Evidence.parent\_evidence} (até três níveis,
\texttt{MAX\_TREE\_DEPTH=3}, \texttt{models.py:318}) por razões de
compatibilidade.

\subsection*{§7.2 --- Suportes de armazenamento}

\textbf{Implementação.} A taxonomia de \texttt{Evidence} (ADR-0010)
inclui os suportes relevantes: \texttt{STORAGE\_MEDIA} (\emph{pen}~USB
ou disco externo), \texttt{INTERNAL\_DRIVE} (HDD/SSD/NVMe),
\texttt{MEMORY\_CARD} (SD/microSD/CF), \texttt{SIM\_CARD},
\texttt{RFID\_NFC\_CARD}, \texttt{DIGITAL\_FILE}. Campos específicos
(IMEI, IMSI, ICCID, MAC, VIN) ficam em
\texttt{Evidence.type\_specific\_data} (\texttt{JSONField}) com
serialização determinística (\texttt{sort\_keys=True},
\texttt{separators=(',', ':')}, \texttt{ensure\_ascii=False},
\texttt{models.py:566--571}).

\subsection*{Auditabilidade}

\textbf{Implementação.} \texttt{AuditLog} (\texttt{models.py:1121--1259})
regista \texttt{user}, \texttt{action} (VIEW, CREATE, EXPORT\_PDF,
EXPORT\_CSV), \texttt{resource\_type}, \texttt{resource\_id},
\texttt{ip\_address}, \texttt{correlation\_id} (UUID por pedido) e
\texttt{timestamp} do servidor. \texttt{IP} é extraído por
\texttt{get\_client\_ip} com \emph{whitelist} de \emph{proxies}
confiáveis (\texttt{audit.py:48--80}). Índices compostos garantem
\emph{queries} rápidas por utilizador, acção, recurso e correlation.

\subsection*{Segurança da infra-estrutura}

\textbf{Implementação.} HTTPS estrito com HSTS de um ano,
\texttt{includeSubDomains} e \texttt{preload}
(\texttt{settings.py:284--286}); CSP \emph{nonce-based} (sem
\texttt{unsafe-inline}) em
\texttt{ContentSecurityPolicyMiddleware}
(\texttt{middleware.py:68--155}); \texttt{Referrer-Policy:
strict-origin-when-cross-origin} (ADR-0007); \emph{rate-limiting}
DRF (anon~10/min, user~60/min, auth~5/min,
\texttt{settings.py:142--150}); paginação obrigatória de 50 itens
(\texttt{BoundedPageNumberPagination}, \texttt{pagination.py}).
Validação externa em
\texttt{docs/compliance/external-tests/}: HSTS \emph{Preload Submission}
aceite, Mozilla HTTP Observatory~A+, Qualys SSL Labs~A+
(18~abr~2026).

\subsection*{Protecção contra IDOR}

\textbf{Implementação.} \texttt{get\_queryset()} de cada \emph{viewset}
filtra automaticamente por \texttt{agent=request.user} quando o perfil
é \texttt{AGENT}, garantindo que um agente só vê os recursos da sua
ocorrência. \texttt{IsOwnerOrReadOnly} valida \emph{ownership} em
escritas. Testes \texttt{AuthorizationIDORTest} em
\texttt{src/backend/core/tests\_api.py} exercitam o caminho cruzado
agente~A vs.\ agente~B, verificando rejeição com \texttt{404}
(em vez de \texttt{403}, para evitar enumeração).

% =============================================================================
\section{Tabela de Cobertura de Requisitos}
% =============================================================================

\begin{longtable}{L{1.6cm} L{8.5cm} L{4cm}}
    \toprule
    \textbf{Req. ISO} & \textbf{Descrição} & \textbf{Estado} \\
    \midrule
    \endhead
    §5.1 & Princípios gerais de manuseamento & Implementado \\
    §5.2 & Identificação de prova digital & Implementado \\
    §5.3 & Recolha e aquisição de prova & Implementado \\
    §5.4 & Preservação de integridade (3 camadas) & Implementado \\
    §5.5 & Cadeia de custódia (7 estados, hash encadeado) & Implementado \\
    §6.1 & Competência de pessoal (RBAC) & Implementado \\
    §6.2 & Documentação do processo (PDF/CSV) & Implementado \\
    §7.1 & Dispositivos digitais & Implementado \\
    §7.2 & Suportes de armazenamento & Implementado \\
    --- & Auditabilidade (\texttt{AuditLog}, \texttt{correlation\_id}) & Implementado \\
    --- & Segurança da infra-estrutura (HSTS preload, CSP nonce, etc.) & Implementado \\
    --- & Protecção contra IDOR & Implementado \\
    \bottomrule
\end{longtable}

% =============================================================================
\section{Notas de Implementação}
% =============================================================================

\begin{itemize}
    \item \textbf{Imutabilidade em três camadas:} mesmo um atacante com
          acesso directo à base de dados não consegue alterar
          retroactivamente uma evidência ou um registo de custódia ---
          \emph{triggers} PostgreSQL bloqueiam \texttt{UPDATE} e
          \texttt{DELETE}.
    \item \textbf{Verificabilidade externa:} qualquer perito
          independente pode recalcular um \emph{hash} de
          \texttt{ChainOfCustody} a partir dos campos públicos do registo
          e do \emph{hash} da raiz da cadeia ($0 \times 64$).
    \item \textbf{Conformidade:} o ForensiQ não é certificado
          ISO/IEC~27037, mas implementa os princípios e os controlos
          previstos na norma. As lacunas para um \emph{deployment} real
          (cifra de \texttt{media/} em repouso, \emph{backup}
          automatizado de \emph{media}, \emph{pentest} externo formal)
          estão registadas em
          \texttt{docs/compliance/rgpd-art32-checklist.md}.
\end{itemize}

% =============================================================================
\section{Histórico de Alterações}
% =============================================================================

\begin{longtable}{L{1.6cm} L{2.6cm} L{10cm}}
    \toprule
    \textbf{Versão} & \textbf{Data} & \textbf{Alteração} \\
    \midrule
    \endhead
    1.0 & 22 mar 2026 & Documento inicial. \\
    1.1 & 12 abr 2026 & Primeira matriz de rastreabilidade. \\
    1.2 & 3 mai 2026  & Refeita a partir do código de produção: estados PT (APREENDIDA \dots DESTRUIDA) em vez de inglês; 18 tipos digital-first (ADR-0010) em vez de USB\_DRIVE/HARD\_DRIVE/SD\_CARD; referências a \texttt{ficheiro:linha} adicionadas. \\
    \bottomrule
\end{longtable}

\end{document}
